% Candidate thesis discussion. Prefer replacing the paper's existing Discussion
% section (baseline lines 312--326); retain its Methods, Results and Limitations.
\section{Chapter discussion}
\label{sec:ch3:chapter-discussion}

\subsection{What the comparison establishes}

This chapter makes expressive production an explicit target for developmental comparison. The main empirical finding is a discrepancy in the joint pattern of vocabulary coverage, lexical diversity, and word-form validity. With larger training corpora, models produce more of the evaluation vocabulary often enough to cross the \meCDI{} threshold. They nevertheless fail to reproduce the accelerated middle portion of the child corpus trajectory. This mismatch is visible across the tested architectures and input registers, although individual measures vary across conditions: audiobook-trained models have especially high lexical diversity, whereas some child-directed-speech conditions approach children's TTR or show substantial declines in non-word production (Figure~\ref{fig:core_metric}). The result is therefore a constraint on the overall production profile, rather than evidence that every model is equally distant from children on every measure.

The temperature experiment strengthens this conclusion by testing one concrete alternative explanation. Among the sampled temperatures, moving toward more diverse output does not yield a configuration that jointly matches the child reference on vocabulary coverage and TTR (Figure~\ref{fig:temp_effect}). It can also increase the rate of forms absent from the reference dictionaries. This is a useful robustness result: the discrepancy is not removed by choosing another value within the tested temperature range. It leaves open whether other elicitation contexts, decoding policies, or training procedures would produce a different relation among the measures.

\subsection{Measurement, benchmark, and model contributions}

The measurement contribution is to specify the observation process through which a vocabulary estimate is obtained. Fixed-size windows make TTR less dependent on the amount of sampled text; production budgets define opportunities to observe words; and the count threshold defines how much observed evidence is required for an item to contribute to the vocabulary score. The benchmark contribution is to combine these decisions with explicit input scales and frequency-matched evaluation sets, allowing a common protocol to be applied to two architectures and two registers. The model comparison then provides empirical constraints on those particular learners under that protocol. These contributions remain useful even when a numerical score cannot be read as an inventory of everything a learner knows.

The distinction matters especially for \meCDI{}. Its threshold was chosen by comparing child corpus counts with a caregiver-report trajectory. This calibration supplies a linking rule between measurements; agreement on the calibration material is not independent validation of lexical knowledge. A target form may remain below threshold because it is weakly represented, because relevant contexts are rarely sampled, or because more probable alternatives occupy the generated output. Conversely, frequent production alone does not establish appropriate reference or contextual use. A stronger validation would apply the fixed calibration to held-out children or corpora and test whether the resulting estimates agree with additional vocabulary measures. The present score is most directly interpretable as the proportion of target forms reaching a production criterion under the specified sampling conditions.

\subsection{Frequency and the meaning of data efficiency}

Frequency analysis identifies where the production discrepancy is largest. The fitted trajectories are more frequency-sensitive for the evaluated models than for the child reference, particularly in the lower-frequency bands (Figure~\ref{fig:freq_effect}). The associated content-word analysis adds an informative distributional difference: the child corpus shows a pronounced rise and subsequent fall in content-word proportion, whereas the model proportions vary within a much narrower range (Figure~\ref{fig:content_prop}). Thus vocabulary coverage and the distribution of output across lexical categories need to be considered together. A model can generate many distinct forms while repeatedly producing too few of the particular content words scored by the benchmark.

This pattern is compatible with an interaction between frequency-weighted predictive training and finite sampling. Under an empirical cross-entropy objective, frequent events contribute repeatedly to the training signal. During generation, low-probability forms also have fewer opportunities to accumulate enough occurrences to pass a fixed count threshold. These considerations motivate the hypothesis that the tested learning and production procedures provide insufficient support for parts of the lexical long tail. The experiment does not isolate the objective as the unique cause: it does not compare cross-entropy with another objective while holding the remaining conditions fixed. Likewise, the content-word proportions do not by themselves identify a causal explanation for TTR. The results locate a frequency-related discrepancy and define hypotheses that can be tested through interventions on memory, objectives, input, or production.

The time axis further limits a numerical claim about data efficiency. Model conditions vary in available corpus size and are trained to convergence, so the amount of distinct text is not the total number of word presentations during optimization. The age-equivalent scale also combines estimated input with increasing output budgets. In the CHILDES conditions, the available training corpus limits further input growth after the 30-month-equivalent checkpoint, while later generation budgets continue to change. Those later observations cannot therefore be attributed wholly to additional learning. Finally, frequency-band estimates that cross the acquisition criterion far beyond the observed range are curve extrapolations, not results from models actually trained on hundreds of months of input. Their magnitude depends on the fitted function and threshold. Within these boundaries, the study shows that increasing corpus size under the tested protocol does not automatically recover the child production profile; it does not determine a universal child--model data multiplier.

\subsection{From a diagnostic discrepancy to a controlled intervention}

The strongest next step is to preserve the production measure while manipulating a candidate source of the discrepancy. An episodic store could test whether access to particular encounters improves low-frequency target coverage; alternative training objectives could test whether changing the weighting of rare events changes that coverage; and grounded or interactive input could test whether referential and conversational context changes which words are produced. Such comparisons should count any additional information and evaluate effects across multiple outcomes, since an improvement in thresholded coverage could coexist with deteriorating lexical diversity or word-form validity.

Chapter~\ref{chp3_chap} takes up one related intervention by asking whether retrieval of stored linguistic instances improves sensitivity to low-frequency syntactic contrasts. It tests the potential contribution of instance-based access using a different outcome and experimental setting. It consequently extends the inquiry into sparse evidence without directly repairing or explaining the expressive-vocabulary curves reported here. Across the two studies, the common question is what information remains available to a predictive learner when relevant evidence is infrequent, and how that availability depends on the task used to measure it.
